\begin{dang}{Hai vân sáng trùng nhau. Tìm số vân}
\begin{itemize}
\item Hiện tượng giao thoa ánh sáng chỉ xảy ra nếu ánh sáng có cùng bước sóng và độ lệch pha không đổi (thông thường người ta tạo ra hai nguồn thứ cấp $S_1$ và $S_2$ cùng pha). 
\item Một nguồn sáng phức tạp gồm nhiều thành phần đơn sắc (bước sóng khác nhau), các thành phần đơn sắc khác nhau không thể giao thoa với nhau. 
\item Khi bài toán đề cập đến việc giao thoa với nhiều bức xạ hay giao thoa của ánh sáng trắng đó là hiện tượng giao thoa của từng bức xạ đơn sắc \bth{\Rightarrow} tạo nên hệ các vân sáng, vân tối riêng biệt chồng chất lên nhau.
\end{itemize}
\paragraph{Giao thoa với hai bức xạ đơn sắc}
\Binhluan{
\begin{itemize}
\item Khi nguồn sáng S phát ra hai bức xạ (ánh sáng) có bước sóng \bth{\lambda_1} và \bth{\lambda_2} thì:
\begin{itemize}
\item Trên màn có hai hệ vân giao thoa ứng với hai ánh sáng có bước sóng \bth{\lambda_1} và \bth{\lambda_2}.
\item Ở vị trí trung tâm có hai vân sáng trung tâm trùng nhau
\item Tại các vị trí M, N,\ldots thì hai vân lại trùng nhau khi 
\begin{align*}
&x_{s1}=x_{s2}\Leftrightarrow k_1i_1=k_2i_2 \Leftrightarrow k_1\lambda_1=k_2\lambda_2\\ 
\Rightarrow &\boder{\dfrac{i_2}{i_1}=\dfrac{\lambda_2}{\lambda_1}=\dfrac{b}{c}}\ \left(\ct{phân số tối giản}\right)
\end{align*}
\end{itemize}
\item Khoảng vân trùng (khoảng cách nhỏ nhất giữa hai vân cùng màu với vân trung tâm)
\begin{align*}
\boder{i_{\equiv}=bi_1=ci_2}\ \ct{hoặc}\ i_{12}=BCNN(i_1,i_2)
\end{align*}
\end{itemize}
}
{Tại O luôn luôn là một vị trí trùng nhau.\\
Về mặt định lượng ta có thể coi như hệ vân giao thoa với một bức xạ "mới" có bước sóng $\lambda_{\equiv}$ và khoảng vân $i_{\equiv}$.
}
%$\bullet$ Vị trí vân sáng trùng nhau của hai hệ vân
%\begin{itemize}
%\item $\dfrac{i_2}{i_1}=\dfrac{\lambda_2}{\lambda_1}=\text{phân số tối giản}=\dfrac{b}{c} \Rightarrow i_{\equiv}=bi_1=ci_2$
%\item Vị trí vân sáng trùng: $x=ni_{\equiv}$ (với n là số nguyên)
%\end{itemize}
%$\bullet$ Số các vị trí vân sáng trùng nhau của hai hệ vân
%\begin{itemize}
%\item Tìm $i_{\equiv}=bi_1=ci_2$
%\item Thay $i_{\equiv}$ vào vị trí của $i$ trong các công thức tính số vân sáng
%\end{itemize}
%$\bullet$ Số vạch sáng, số vạch sáng đơn sắc:Gọi $N_1$, $N_2$ lần lượt là số vân sáng khi giao thoa lần lượt với $\lambda_1$ và $\lambda_2$; $N_{\equiv}$ là số vân sáng trùng nhau của hai bức xạ.\\
%\textbf{Bài toán 1:} Tìm số vân sáng quan sát được
%\begin{itemize}
%\item Mỗi vân sáng là một vạch sáng, nhưng nếu vân sáng hệ này trùng vân sáng hệ kia chỉ cho ta một vạch sáng (vân sáng trùng).
%\item Tìm số vân sáng của hệ 1 là $N_1$ và của hệ 2 là $N_2$.
%\item Tìm số vân sáng trùng là $N_{\equiv}$.
%\item Số vân sáng là: $N_{vs}=N_1+N_2-N_{\equiv}$.
%\end{itemize}
%\textbf{Bài toán 2:} Tìm số vạch sáng nằm giữa vân sáng bậc $k_1$ của $\lambda_1$ và vân sáng bậc $k_2$ của $\lambda_2$
%\begin{itemize}
%\item \textit{Cách 1}: Dùng hình vẽ
%\begin{itemize}
%\item Vẽ các vân trùng cho đến bậc $k_1$ của hệ 1 và bậc $k_2$ của hệ 2.
%\item Từ hình vẽ ta xác định được số vạch sáng.
%\end{itemize}
%\item \textit{Cách 2}: Dùng toạ độ
%\begin{itemize}
%\item Số vân sáng của $\lambda_1$: $k_1i_1 \leq m_1i_1 \leq k_2i_2 \Rightarrow N_1$
%\item Số vân sáng của $\lambda_2$: $k_1i_1 \leq m_2i_2 \leq k_2i_2 \Rightarrow N_2$
%\item Số vân sáng của $\lambda_1 \equiv \lambda_2$ \ (hay $\lambda_{\equiv}$): $k_1i_1 \leq m_{\equiv}i_{\equiv} \leq k_2i_2 \Rightarrow N_{\equiv}$
%\item Số vân sáng là: $N_{vs}=N_1+N_2-N_{\equiv}$.
%\end{itemize}
%\end{itemize}
%\textbf{Bài toán 3:} Tìm số vân đơn sắc
%\begin{itemize}
%\item Vân trùng không phải vân đơn sắc
%\item Số vân sáng đơn sắc: $N_{\text{đ}s}=N_1+N_2-2N_{\equiv}$
%\item Số vân sáng đơn sắc của $\lambda_1$: $N_{\lambda1}=N_1-N_{\equiv}$
%\item Số vân sáng đơn sắc của $\lambda_2$: $N_{\lambda2}=N_2-N_{\equiv}$
%\end{itemize}
%\begin{note}
%Để tìm $N_1$ và $N_2$ ta chú ý kiến thức đã học ở bài giao thoa với ánh sáng đơn sắc
%\end{note}
\end{dang}
